\section*{Capítulo \ref{ch:g8:puberty} --- A puberdade}

\begin{solution}{exo:g8:puberty:1}
A etapa em que o corpo de uma criança vira o corpo de um adulto --- com
os órgãos reprodutores amadurecendo e começando a produzir \omterm{def:g8:sexual-reproduction:gametes}{gametas}, e
a forma adulta assumida. Ela é o cerne biológico da adolescência.
\end{solution}

\begin{solution}{exo:g8:puberty:2}
Meninas: as mamas se desenvolvem, os quadris se alargam, vêm as
primeiras menstruações (com os pelos do corpo). Meninos: a voz
engrossa, os \omterm{def:g1:my-body:joint}{ombros} se alargam, começa a produção de \omterm{def:g4:animal-reproduction:sexual}{espermatozoides}
(com os pelos do rosto). Nos dois: o salto de crescimento, a \omterm{def:g1:the-five-senses:sense}{pele} mais
oleosa, o \omterm{rem:g7:removing-wastes:sweat}{suor} mudado, o tempo instável de humor.
\end{solution}

\begin{solution}{exo:g8:puberty:3}
Nas meninas, as primeiras menstruações --- o início dos ciclos; nos
meninos, os testículos começando a produzir \omterm{def:g4:animal-reproduction:sexual}{espermatozoides}.
\end{solution}

\begin{solution}{exo:g8:puberty:4}
Um mensageiro químico liberado no \omterm{prop:g4:journey-of-food:absorb}{sangue} por um órgão especializado,
levado a toda parte pela circulação e obedecido pelos órgãos
construídos para responder a ele.
\end{solution}

\begin{solution}{exo:g8:puberty:5}
Ela começa num centro de comando na base do \omterm{def:g5:brain-in-command:system}{cérebro}, cujos \omterm{def:g8:puberty:hormones}{hormônios}
acordam os órgãos reprodutores; os \omterm{def:g8:puberty:hormones}{hormônios} deles, por sua vez,
comandam a reconstrução do corpo inteiro --- então ela termina em toda
parte aonde o \omterm{prop:g4:journey-of-food:absorb}{sangue} chega.
\end{solution}

\begin{solution}{exo:g8:puberty:6}
Os dois são normais: o calendário é pessoal --- cedo e tarde são os
dois saudáveis, a ordem das mudanças é variável e o destino é
compartilhado. Estar a anos de distância aos treze não diz quase nada
sobre nenhum dos dois.
\end{solution}

\begin{solution}{exo:g8:puberty:7}
Porque o mensageiro viaja no \omterm{prop:g4:journey-of-food:absorb}{sangue}, que visita cada órgão de qualquer
jeito: uma transmissão só chega à laringe, aos \omterm{def:g1:my-body:joint}{ombros} e à \omterm{def:g1:the-five-senses:sense}{pele} por
igual, sem fiação nenhuma. Os órgãos que respondem são os que foram
construídos para escutar aquele mensageiro --- endereçamento por
receptor, e não por rota.
\end{solution}

\begin{solution}{exo:g8:puberty:8}
Causa um: os próprios \omterm{def:g8:puberty:hormones}{hormônios} da \omterm{def:g8:puberty:puberty}{puberdade} chegam ao \omterm{def:g5:brain-in-command:system}{cérebro}. Causa
dois: o \omterm{def:g5:brain-in-command:system}{cérebro} está ao mesmo tempo se religando rumo à forma adulta
dele. Alívio: é tempo instável --- normal, explicável e passageiro.
\end{solution}

\begin{solution}{exo:g8:puberty:9}
A matéria de construção acima de tudo --- o salto constrói \omterm{def:g3:bones-and-muscles:skeleton}{osso} e
\omterm{def:g3:bones-and-muscles:muscle}{músculo} --- com quem constrói \omterm{def:g3:bones-and-muscles:skeleton}{osso} vindo dos laticínios e a \omterm{prop:g3:balanced-diet:jobs}{energia} que
dura vindo da família dos alimentos com amido: o raciocínio da dieta
radical do \cref{exo:g7:digestion-and-nutrients:11}, rodado para a
frente.
\end{solution}

\begin{solution}{exo:g8:puberty:10}
Um canteiro de obras é mais vulnerável no meio da construção: o
tabaco, o álcool e os demais atingem órgãos --- o \omterm{def:g5:brain-in-command:system}{cérebro} em primeiro
lugar --- cuja fiação e cujos tecidos ainda estão sendo assentados.
Estrago numa estrutura pronta é reparo; estrago numa estrutura pela
metade é mudança de projeto.
\end{solution}

\begin{solution}{exo:g8:puberty:11}
Confiáveis: o enfermeiro da escola, a médica de família (formados,
sigilosos, impossíveis de escandalizar), este livro (feito para ser
relido e conferido). Não confiáveis: o vídeo viral (ele vende,
negociando com a incerteza), a lenda do pátio (palpites recontados,
sem endereço para correções).
\end{solution}

\begin{solution}{exo:g8:puberty:12}
Alcance: um mensageiro muda a voz, a \omterm{def:g1:the-five-senses:sense}{pele}, o \omterm{def:g3:bones-and-muscles:skeleton}{esqueleto} e o humor ao
mesmo tempo --- o \omterm{prop:g4:journey-of-food:absorb}{sangue} entrega em toda parte (observação: a
reconstrução coordenada). Coordenação: as mudanças seguem um
cronograma só, embora \omterm{def:g5:brain-in-command:system}{nervo} nenhum ligue os canteiros (observação: a
barba e a voz chegando juntas). Diferença em relação aos \omterm{def:g5:brain-in-command:system}{nervos}: sem
fiação e sem velocidade de fração de segundo --- as ordens levam meses e
agem no corpo inteiro, enquanto as mensagens da
\cref{prop:g5:brain-in-command:loop} levam milésimos de segundo e um
caminho só (observação: os anos lentos da \omterm{def:g8:puberty:puberty}{puberdade} contra o quinto de
segundo de uma pegada).
\end{solution}

\begin{solution}{pb:g8:puberty:1}
\textbf{1.} As quatro mostram o salto da adolescência. L: salto quase
completo, achatando rumo à altura adulta. M: salto ainda não
começado --- ainda na subida constante da infância. N: no meio do
salto. O: um salto cedo, dos nove aos doze, agora chegando ao platô.

\textbf{2.} ``A sua curva é exatamente uma curva normal antes do
salto: o calendário é pessoal, chegar tarde é tão saudável quanto
chegar cedo, e a sua subida constante diz que o salto está vindo ---
curvas como a sua ficam íngremes em um ou dois anos e chegam às mesmas
alturas adultas de todo mundo.''

\textbf{3.} O platô encerra o salto de \emph{altura} --- o crescimento
em estatura. O resto do trabalho da \omterm{def:g8:puberty:puberty}{puberdade} --- o corpo se
preenchendo, o amadurecimento, a reforma do \omterm{def:g5:brain-in-command:system}{cérebro} --- continua depois
dele; a O está adiantada, e não terminada.

\textbf{4.} Porque a altura mistura herança e idade --- uma criança
alta antes do salto pode medir mais que outra baixa no meio do salto.
A inclinação mostra o processo: íngreme quer dizer salto, achatando
quer dizer fim do salto, seja qual for a altura da linha.

\textbf{5.} O despertar do centro de comando na base do \omterm{def:g5:brain-in-command:system}{cérebro}; os
\omterm{def:g8:puberty:hormones}{hormônios} dele ordenaram que os órgãos reprodutores acordassem, e os
\omterm{def:g8:puberty:hormones}{hormônios} deles --- transmitidos pelo \omterm{prop:g4:journey-of-food:absorb}{sangue} --- ordenaram as placas de
crescimento e todo o resto.

\textbf{6.} \omterm{prop:g1:healthy-habits:needs}{Sono}: o crescimento e a reforma do \omterm{def:g5:brain-in-command:system}{cérebro} acontecem de
noite --- noites curtas matam de fome os dois. Alimentação: o salto é
pago em matéria --- a matéria de construção e as famílias do prato o
financiam. (O movimento e o nada de sabotagem completam a lista de
cuidados.)

\textbf{7.} Os \omterm{def:g8:puberty:hormones}{hormônios} chegam ao \omterm{def:g5:brain-in-command:system}{cérebro} do N, e o \omterm{def:g5:brain-in-command:system}{cérebro} está no
meio da religação --- química mais obra, normal e passageiro. Na
prática: o relógio escorrega para tarde, mas a escola não; o N precisa
defender as horas com um encerramento do dia mais cedo e sem \omterm{def:g6:exploring-our-environment:conditions}{luz} de
tela (\cref{prop:g5:healthy-choices:screens}).

\textbf{8.} Porque o consultório é formado, sigiloso e impossível de
escandalizar --- as perguntas que a curva não mostra também têm um
endereço certo, e a lenda e os cantos da internet não são esse
endereço.

\textbf{9.} L: salto terminando, normal --- proteger o \omterm{prop:g1:healthy-habits:needs}{sono} nos anos de
prova. M: antes do salto, normal --- paciência e prato. N: no meio do
salto, normal --- horas de \omterm{prop:g1:healthy-habits:needs}{sono} e matéria de construção. O: salto cedo
concluído, normal --- movimento para os anos de preenchimento do corpo.

\textbf{10.} Porque normal é uma estrada larga, e não uma linha:
quatro calendários saudáveis, a anos de distância, todos chegando.

\textbf{11.} Nada: a maquinaria está íntegra e o calendário dela é
pessoal --- não há defeito para consertar nem alavanca que valha a \omterm{def:g1:animals-around-us:covering}{pena}
puxar numa curva saudável; a receita honesta é tempo (com a lista de
cuidados enquanto isso passa).
\textbf{12.} ``A \omterm{def:g8:puberty:puberty}{puberdade} é a reconstrução do corpo, de criança para
adulto, ordenada por \omterm{def:g8:puberty:hormones}{hormônios} --- comandada por mensageiros no
\omterm{prop:g4:journey-of-food:absorb}{sangue}, num calendário pessoal, e financiada por \omterm{prop:g1:healthy-habits:needs}{sono}, comida e
movimento.''
\end{solution}
